\chapter*{集合专题练习}
\begin{enumerate}
    \item 集合的概念
    \begin{enumerate}
        \item 用符号$\in , \notin$填空： \\
        \begin{enumerate*}
            \item $0$\blkda{  $\in$  }$\mathbf{N}$
            \item $0 $\blkda{  $\notin$  }$\mathbf{N^*}$
            \item $0 $\blkda{  $\notin$  }$\varnothing$
            \item $\sqrt{5} $\blkda{  $\notin$  }$\mathbf{N}$
            \item $-\sqrt{9}$\blkda{  $\in$  }$\mathbf{Z^-}$ \\
            \item $-\dfrac{1}{2} $\blkda{  $\in$  }$ \mathbf{Q}$
            \item $\pi $\blkda{  $\notin$  }$ \mathbf{Q}$
            \item $3.14 $\blkda{  $\in$  }$ \mathbf{R^+}$ 
            \item $0$\blkda{  $\notin$  }$\{x | x^2 + 1 = 0\}$ \\
            \item $\sqrt{2} $\blkda{  $\notin$  } $\{x | x \geq \sqrt[3]{3}$
            \item $(0, 0) $\blkda{  $\in$  }$  \{(x,y) | x + y - 1 \leq 0\} $
        \end{enumerate*}
        \item 用列举法表示集合：
        \begin{enumerate}
            \item $\{x|3x-1 \leq 11, x \in \mathbf{N}\}$
            \item $\{x|x \in N \mbox{且} \dfrac{6}{x+2} \in \mathbf{N}\}$
            \item $\{(x, y) | x^2 + y^2 \leq 2, x \in \mathbf{Z}, y \in \mathbf{Z}\}$
        \end{enumerate}
        \begin{jieda}
            \begin{enumerate}
                \item $\{0,1,2,3,4\}$
                \item $\{0,1,4\}$
                \item $\{(0,0), (1,1), (-1,1), (1, -1), (-1,-1), (0, 1), (0, -1), (1, 0), (-1, 0)\}$
            \end{enumerate}
        \end{jieda}
        \item 用描述法表示下列集合：
        \begin{enumerate}
            \item 被3除余1的所有自然数组成的集合；
            \item 比1大又比10小的所有实数组成的集合；
            \item 平面直角坐标系中坐标轴上所有点组成的集合.
        \end{enumerate}
        \begin{jieda}
            \begin{enumerate}
                \item $\{x | x = 3n+1, n \in \mathbf{N} \}$
                \item $\{x |1 < x < 10 \}$
                \item $\{(x, y) | xy = 0\}$
            \end{enumerate}
        \end{jieda}
        \item 用适当的符合填空：
        \begin{enumerate}
            \item  $\varnothing$ \blkda{=} $ \{x \in \mathbf{R} | x^2 + x + 1 = 0\}$
            \item $\{x | x = 4m -1, m \in \mathbf{Z}\}$ \blkda{$\subset$}  $\{x | x = 2n + 1, n \in \mathbf{Z}\}$
            \item $\{y | |x| + |y| = 1\}$  \blkda{$=$}  $\{y | x^2 + y^2 \leq 1\}$
            \item $\{y | y = x^2-4x+3 \}$   \blkda{$=$}   $\{y | y = x^2-1\}$
        \end{enumerate}
    \end{enumerate}
    \item 判断、证明集合间的关系
    \begin{enumerate}
        \item 已知集合$M = \{x | x = m + \dfrac{1}{6}, m \in \mathbf{Z}\}, N = \{x | x = \dfrac{n}{2} - \dfrac{1}{3}, n \in \mathbf{Z}\}, P = \{x | x = \dfrac{p}{2} + \dfrac{1}{6}, p \in \mathbf{Z}\}$，试判断$M, N, P$三者的关系，并说明理由.
        \begin{jieda}
            首先进行变形有$M = \{x | x = \dfrac{6m+1}{6}, m \in \mathbf{Z}\}, N = \{x | x = \dfrac{3n-2}{6}, n \in \mathbf{Z}\}, P = \{x | x = \dfrac{3p+1}{6}, p \in \mathbf{Z}\}$ \\
            $\forall x \in N, x = \dfrac{3n-2}{6} = \dfrac{3(n-1)+1}{6}$，$n-1 \in \mathbf{Z}$，故$x \in P$，因此$N \subseteq P$。\\
            $\forall x \in P, x = \dfrac{3p+1}{6} = \dfrac{3(p+1)-2}{6}$，$p+1 \in \mathbf{Z}$，故$x \in N$，因此$P \subseteq N$。\\
            因此$N = P$ \\
            $\forall x \in M, x = \dfrac{6m+1}{6} = \dfrac{3(2m)-2}{6}$，$2m \in \mathbf{Z}$，故$x \in P$，因此$M \subseteq P$。\\
            而$p = 1$时$\dfrac{2}{3} \in P$，但是$\dfrac{2}{3} \notin M$，因此$P \neq M$ \\
            因此$M \subset P$
        \end{jieda}
        \item 设集合 $M=\{x|x=\dfrac{k}{2} + \dfrac{1}{4}, k \in \mathbf{Z}\}$, $N=\{x|x=\dfrac{k}{4} + \dfrac{1}{2}, k \in \mathbf{Z}\}$, 则\dotfill{(\kh{B})}. 
        \xx{$M=N$}{$M \subset N$}{$M \supset N$}{$M \cap N = \varnothing $}
        \begin{jieda}
            $M=\{x|x=\dfrac{2m+1}{4}, m \in \mathbf{Z}\}$， $N=\{x|x=\dfrac{n+2}{4}, n \in \mathbf{Z}\}$\\
            对任意的$m$，可以取$n = 2m-1$，则$\dfrac{n+2}{4} = \dfrac{2m+1}{4}$，因此$M$中的元素均在$N$中。\\
            但取$n$为偶数时，例如$n = 0$，不存在$m$使得$\dfrac{n+2}{4} = \dfrac{2m+1}{4}$，因此$N$中有部分元素不在$M$中。故$M \subset N$\\
            \ps 带入特殊的$k$值，1，2，3等可以看出规律
        \end{jieda}
        \item 设集合 $M=\{x|x=\dfrac{k}{4} - \dfrac{1}{8}, k \in \mathbf{Z}\}$, $N=\{x|x=\dfrac{k}{2} \pm \dfrac{3}{8}, k \in \mathbf{Z}\}$, 则\dotfill{(\kh{A})}
        \xx{$M=N$}{$M \subset N$}{$M \supset N$ }{$M \cap N = \varnothing $ }
        \begin{jieda}
            $M=\{x|x=\dfrac{2m-1}{8}, m \in \mathbf{Z}\}$， $N=\{x|x=\dfrac{4n \pm 3}{8}, n \in \mathbf{Z}\}$
            \begin{enumerate}
                \item 对任意的$m$，若$m$为奇数，则可以取$n = \dfrac{m+1}{2}$使得$\dfrac{4n-3}{8} = \dfrac{2m-1}{8}$，若$m$为偶数，则可以取$n = \dfrac{m}{2}-1$使得$\dfrac{4n+3}{8} = \dfrac{2m-1}{8}$，
                \item 对任意的$n$，可以取$m_1 = 2n+2, m_2 = 2n-2$，则$\dfrac{4n+3}{8} = \dfrac{2m_1-1}{8}, \dfrac{4n-3}{8} = \dfrac{2m_2-1}{8}$，因此$N$中的元素均在$M$中。
            \end{enumerate}
            故$M = N$
        \end{jieda}
    \end{enumerate}
    \item 集合的运算
    \begin{enumerate}
        \item 已知全集$U = \{1,2,3,4,5\}$，集合$M = \{1,2\}, N = \{3,4\}$，则$\overline{M \cup N} =$\blkda{$\{5\}$}
        \begin{jieda}
            根据并集的运算可知$M\cup N = \{1,2,3,4\}$，故其补集$\overline{M\cup N} = \{5\}$ 
        \end{jieda}
        \item 设全集$U = \{-2, -1, 0, 1, 2\}$，集合$A = \{0,1,2\}, B = \{-1, 2\}$，则$A \cap \overline{B} = $\blkda{$\{0,1\}$}
        \begin{jieda}
            首先得到$\overline{B} = \{-2,0,1\}$故$A \cap \overline{B} = \{0,1\}$
        \end{jieda}
        \item 已知集合$U = \{1,2,4,6,8\}, M = \{x | x^2-3x+2=0\}, N = \{x | x = 4a, a \in M\}$，则$\overline{M \cup N} = $ \blkda{$\{6\}$}
        \begin{jieda}
            $M = \{1,2\}, N = \{4, 8\}$，故$\overline{M \cup N} = \{6\}$
        \end{jieda}
        \item 已知集合$M = \{x | x^2 + x -2 \leq 0\}, Q = \{x \in \mathbf{N} | |x|\ \leq 2\}$，则$M \cap Q = $ \blkda{$\{0, 1\}$}
        \begin{jieda}
            首先将$M$转化成区间表示：$M = [-2, 1]$，将$Q$转化为列举法表示$Q = \{0,1,2\}$，故$M \cap Q = \{0, 1\}$
        \end{jieda}
        \item 已知集合$A = \{x \in \mathbf{Z} | x^2 - 8x + 15 \leq 0\}, B = \{x | x < 5\}$，则$A \cap B = $\blkda{$\{3,4\}$}
        \begin{jieda}
            $A = \{3,4,5\}, B = (-\infty, 5)$，故$A \cap B = \{3,4\}$
        \end{jieda}
        \item 已知$P = [1,3], M = \{x | x^2 \geq 4\}$，则$P \cup \overline{M} = $\blkda{$(-2,3]$}
        \begin{jieda}
            $\overline{M} = \{x | x^2 < 4\}$，即$M = (-2, 2)$。故$P \cup \overline{M} = (-2,3]$ 
        \end{jieda}
        \item 已知全集$U = \{x | x \leq 4\}$，集合$A = \{x | -2 < x < 3\}, B = \{x | -3 \leq x \leq 2\}$，求$A \cap B, \quad \overline{A}\cup B, \quad A\cap\overline{B}$
        \begin{jieda}
            $A = (-2,3), B = [-3,2], \overline{A} = (-\infty, -2] \cup [3, 4], \overline{B} = (-\infty, -3) \cup (2, 4]$ \\
            故$A \cap B = (-2, 2], \quad \overline{A}\cup B = (-\infty, -2] \cup [3, 4],\quad  A\cap\overline{B} = (2, 3)$
        \end{jieda}
    \end{enumerate}
    \item 描述法与表示符号
    \begin{enumerate}
        \item 求下列集合运算的结果
            \begin{enumerate}
                \item 已知集合$A = \{y | y = x^2+1\}$，$B = \{y | y = \dfrac{1}{x}\}$，求$A \cap B$，$A \cup B$
                \item 已知集合$A = \{x | y = x^2+1\}$，$B = \{x | y = \dfrac{1}{x}\}$，求$A \cap B$，$A \cup B$
                \item 已知集合$A = \{(x, y) | y = x^2+1\}$，$B = \{(x, y) | y = \dfrac{x^2-1}{x-1}\}$，求$A \cap B$
            \end{enumerate}
            \begin{jieda}
                \begin{enumerate}
                    \item $A = [1, +\infty)$，$B = (-\infty, 0) \cup (0, +\infty)$，故$A \cap B = [1, +\infty)$，$A \cup B = (-\infty, 0) \cup (0, +\infty)$
                    \item $A = \mathbf{R}$，$B = (-\infty, 0) \cup (0, +\infty)$，故$A \cap B = (-\infty, 0) \cup (0, +\infty)$，$A \cup B = \mathbf{R}$
                    \item 集合$A$是抛物线$y = x^2+1$上的点，集合$B$是曲线$y = \dfrac{x^2-1}{x-1}$上的点，即直线$y = x+1$上除$(1,2)$之外的点。故两者交点仅有$(0, 1)$，因此$A \cap B$为$\{(0,1)\}$。
                \end{enumerate}
            \end{jieda}
        \item 设集合$M = \{y | y = \sqrt{x} - 1 \}， N = \{y | y = -x^2+1 \}$，则$M \cap N = $\blkda{$[-1, 1]$}
        \begin{jieda}
            $M = [-1, +\infty), N = (-\infty, 1]$
        \end{jieda}      
        \item 设集合$M = \{x | y = \sqrt{x} - 1 \}， N = \{y | y = -x^2+1 \}$，则$M \cap N = $\blkda{$[0, 1]$}
        \begin{jieda}
            $M = [0, +\infty), N = (-\infty, 1]$ \\
            \ps 描述法需要看清表示元素的符号。
        \end{jieda}
        \item 设集合$M = \{(x, y) | y = \sqrt{x} - 1 \}， N = \{(x, y) | y = -x^2+1 \}$，则$M \cap N = $\blkda{$\{(1, 0)\}$}
        \begin{jieda}
            $\sqrt{x} - 1 = -x^2+1$，解得$x = 1$，故$M \cap N = \{(1, 0)\}$ \\
            \ps 本题两个集合为两曲线，求交集即两曲线求交点。两曲线求交点联立！数集不联立！
        \end{jieda}
    \end{enumerate}
    \item 集合的运算与集合间关系的转化
    \begin{enumerate}
        \item 集合$M = \{x | ax + 1 = 0\}$，$N = \{x | x^2-3x+2=0\}$， 且$M \cap N = M$，则 $a = $ \blkda[3]{0, -1, 或 -$\dfrac{1}{2}$}.
        \begin{jieda}
            $N = \{1, 2\}$. 
            \begin{dingautolist}{172}
                \item 若$a = 0$, 则$M = \varnothing$, 符合。
                \item 若$a \neq 0$，$M = \left\{-\dfrac{1}{a} \right\}$， 因此 $\dfrac{1}{a} = -1 \mbox{或} -2$, 即$a = -1 \mbox{或} -\dfrac{1}{2}$
            \end{dingautolist}
            综上，$a = 0, -1$, 或 $-\dfrac{1}{2}$
        \end{jieda}
        \item 已知集合$A = \{x | -2 \leq x \leq 7\}, B = \{x | m+1 < x < 2m-1\}$，若$A \cup B = A$，则$m$的取值范围为\blkda{$(-\infty, 4]$}
        \begin{jieda}
            首先根据$A \cup B = A$可知$B \subset A$。由于$B$是描述法表示的集合，因此需要分其是否为空集进行讨论。
            \begin{enumerate}
                \item $B = \varnothing$，此时$m+1 \geq 2m -1$，即$m \leq 2$，满足题意
                \item $B \neq \varnothing$，此时应有$m+1 \geq -2, 2m-1 \leq 7, m > 2$，故$m \in (2, 4]$
            \end{enumerate}
            综上，$m \in (-\infty, 4]$
        \end{jieda}
    \end{enumerate}
    \item 元素个数与子集个数的转化
    \begin{enumerate}
        \item 已知集合$A$中有10个元素，那么集合$A$的非空真子集个数为\blkda{$1022$}
        \item 若集合$S = \{x | \dfrac{6}{x-2}\ \in \mathbf{Z} \mbox{且} x \in \mathbf{Z} \}$, 则集合S的非空真子集个数为\blkda{$254$}.
        \begin{jieda}
            由于 $\dfrac{6}{x-2}\ \in \mathbf{Z}$, 因此$x-2 \in \{-6, -3, -2, -1, 1, 2, 3, 6\}$，即 $S = \{-4, -1, 0, 1, 3, 4, 5, 8\}$， 因此集合$S$的元素个数为8， 因此其非空真子集的个数为$2^8-2=254$ \\
            \ps 含有n个元素的有限集的非空真子集个数为$2^n -2$
        \end{jieda}
    \end{enumerate}
        
    \item 有限集推理求值类问题
    \begin{enumerate}
        \item 设集合$A = \{1, 4, x\}$，$B = \{1, x^2\}$，且$A \cap B = B$，则$x = $\blkda[3]{$0 \mbox{或} -2 \mbox{或} 2$}
        \begin{jieda}
            首先对条件进行转换，$A \cap B = B$即$B \subseteq A$，根据集合中元素的无序性可知需要进行分类讨论。
            \begin{enumerate}
                \item $x^2 = x$，此时对应于$x = 0 \mbox{或} 1$，但是经过检验发现$x = 1$使得集合中元素违背互异性。
                \item $x^2 = 4$，此时对应于$x = 2 \mbox{或} -2$，均不违背互异性。
            \end{enumerate}
            综上，$x = 0 \mbox{或} -2 \mbox{或} 2$
        \end{jieda}
        \item 设集合$A = \{a, a^2, ab\}$，集合$B = \{1, a, b\}$，若$A = B$，求实数$a,b$
        \begin{jieda}
            首先可以判断出来元素1应该是突破口。在集合A中，$a \neq 1$否则违背互异性。由于集合元素的无序性，此时要分类讨论。
            \begin{enumerate}
                \item $a^2 = 1$，此时对应于$a = -1, ab = b$，所以$b = 0$
                \item $ab = 1$，此时对应于$a = \dfrac{1}{b}, a^2 = b$，解得$a = b = 1$（舍去）
            \end{enumerate}
            综上，$a = -1, b = 0$
        \end{jieda}
        \item 设集合 $A = \left\{  {{a}_{1},{a}_{2},{a}_{3},{a}_{4}}\right\}$ ,若集合 ${A}$ 中所有三个元素的子集中的三个元素之和组成的集合为 $B = \{ 2,3,4,6\}$ ,则集合 $A =$ \blkda{$\{-1,1,2,3\}$}
        \begin{jieda}
            设$s = a_1 + a_2 + a_3 + a_4$，则$B = \{s-a_1, s-a_2, s-a_3, s-a_4\}$，求$B$中元素和可知$4s-s = 15$，故$s = 5$，因此$A = \{-1,1,2,3\}$
        \end{jieda}
        \item 已知 $A = \left\{  {{a}_{1},{a}_{2},{a}_{3},{a}_{4}}\right\}  ,B = \left\{  {{a}_{1}^{2},{a}_{2}^{2},{a}_{4}^{2}}\right\}$ 且 ${a}_{1} < {a}_{2} < {a}_{3} < {a}_{4}$ ,其中 ${\mathbf{a}}_{\mathbf{i}} \in  \mathbf{Z}(i = 1,2,3,4)$,若 $A \cap  B = \left\{  {{a}_{2},{a}_{3}}\right\}  ,{a}_{1} + {a}_{3} = 0$ ,且 $A \cup  B$ 的所有元素之和为 56 ,求 ${a}_{3} + {a}_{4} =$ \blkda{8}.
        \begin{jieda}
            由 ${a}_{1} + {a}_{3} = 0$ 得 ${a}_{1} =  - {a}_{3}$ ,所以 ${a}_{1}^{2} = {a}_{3}^{2}$，故$a_2^2 < a_1^2 = a_3^2 < a_4^2$
    
            由${\mathbf{a}}_{\mathbf{i}} \in  \mathbf{Z}(i = 1,2,3,4)$可知$a_2^2 \geq 0, a_3^2 \geq 1, a_4^2 \geq 4$，因此必有$a_4^2 > a_4$，故$\left\{\begin{array}{l}
                 a_2^2 = a_2 \\
                 a_3^2 = a_3
            \end{array}\right.$，即$\left\{\begin{array}{l}
                 a_2 = 0 \\
                 a_3 = 1
            \end{array}\right.$，故$A \cup  B = \{-1, 0, 1, a_4, a_4^2\}$，故$a_4 = 7$，因此$a_3+a_4 = 8$
        \end{jieda}
        \item *设集合$A=\{1, 2, m\}$，其中$m$为实数，令$B=\{a^2|a \in A\}$， $C = A \cup B$。若$C$中所有元素之和为6，则$C$中所有元素之积为\blkda{$-8$}
        \begin{jieda}
            $B = \{1, 4, m^2\}$，则$1, 2, 4, m, m^2 \in C$。根据$C$中元素和为6，而$1+2+4 > 6$，又因为$m^2+m > -1$恒成立，故应有$m < 0$ 且 $m^2 \in \{1, 2, 4\}$。故可知$m = -1$, $C = \{1, 2, 4, -1\}$，因此$C$中所有元素之积为-8.
        \end{jieda}
    \end{enumerate}
    \item 根据集合间关系求值
    \begin{enumerate}
        \item 已知集合$A = \{x | a-2 \leq x \leq a+2\}, B = \{x|x \leq 2 \mbox{或} x \geq 4\}$，若$A \subseteq B$，则$a$的取值范围是\blkda[4]{$(-\infty, 0] \cup [6, +\infty)$}.
        \begin{jieda}
            $A \subseteq B$即$a+2 \leq 2$或$a-2 \geq 4$，解得$a \in (-\infty, 0] \cup [6, +\infty)$
        \end{jieda}
        \item 设$a, b \in \mathbf{R}$，集合$A = \{x | x^2 - 3x - 4 \leq 0\}, B = \{x | x^2 - (2a+1)x < 0\}, C = [2b-1, b+5]$.
        \begin{enumerate}
            \item 若$B \subseteq A$，求$a$的取值范围.
            \item 若$A \subseteq C$，求$b$的取值范围.
            \item 若$a = 6, C \subseteq B$，求$b$的取值范围.
        \end{enumerate}
        \begin{jieda}
            $A = [-1, 4]$ \\
            若$a = -\dfrac{1}{2}$，则$B = \varnothing$，
            若$a > -\dfrac{1}{2}$，则$B = (0, 2a+1)$，
            若$a < -\dfrac{1}{2}$，则$B = (2a+1, 0)$，
            \begin{enumerate}
                \item $B = \varnothing$满足要求，若$B \neq \varnothing$，则有$2a+1 \in [-1, 4]$，因此有$a \in [-1, \dfrac{3}{2}]$
                \item $A \subseteq C$即$2b-1 \leq -1 < 4 \leq b+5$，故$b \in [-1, 0]$
                \item $a = 6, B = (0, 13)$，$C \subseteq B$即$0 < 2b-1 < b + 5 < 13$，故$b \in (\dfrac{1}{2}, 6)$
            \end{enumerate}
        \end{jieda}
    \end{enumerate}
    \item 转化的等价性
    \begin{enumerate}
        \item 若集合$A = \{x | (a-1)x^2+3x-2 = 0\}$仅有一个真子集，则实数$a$的值为\blkda[3]{$a = 1 \mbox{或} -\dfrac{1}{8}$}
        \begin{jieda}
            由于只有一个真子集，故集合中只有一个元素。因为二次项系数有参数$a$，所以需要分类讨论
            \begin{enumerate}
                \item $a = 1$，此时方程退化回一次方程，只有一个实根，满足题意  
                \item $a \neq 1$，此时方程为二次方程，若只有一个实根，则判别式为0，即$9 + 8(a-1) = 0$，得$a = -\dfrac{1}{8}$
            \end{enumerate}
        \end{jieda}
        \item 设全集$U = \{(x, y) | x, y \in \mathbf{R}\}$, $A = \{(x, y) | \dfrac{y-3}{x-2}=1, x, y \in \mathbf{R}, 且x \neq 2\}$, $B = \{(x, y) | y=x+1, x, y \in \mathbf{R}\}$, 则$\overline{A}\cap B = $ \dotfill{(\kh{C})}.
        \xx{$\overline{A}$}{$\overline{B}$}{$\{(2, 3)\}$}{$\varnothing$}
        \begin{jieda}
            $A = \{(x, y) | y=x+1, x, y \in \mathbf{R}, 且x \neq 2\}$, 因此$A \subsetneqq B$。$(2, 3) \in B \mbox{且} (2, 3) \notin A$，故选C。
        \end{jieda}
        \item 已知集合$A = \{x | x^2 - ax + a^2 - 19 = 0\}, B = \{x|x^2-5x+6=0\}, C = \{x | x^2 + 2x - 8 = 0\}$，且满足$A \cap B \neq \varnothing, A \cap C = \varnothing$，求参数$a$的值.
        \begin{jieda}
            首先可以解方程求出$B = \{2,3\}, C = \{-4,2\}$。由$A \cap B \neq \varnothing, A \cap C = \varnothing$可知$2 \notin A, 3 \in A$; 代入$A$中有$9 - 3a + a^2 - 19 = 0$，解得$a = 5 \mbox{或} -2$; 代入$A$中检验发现当$a = 5$时$A = B, 2 \in A$不符合题意。故$a = -2$.
        \end{jieda}
    \end{enumerate}
\end{enumerate}
